Table~\ref{tab:stage1-main-results} is the core Stage 1 comparison.
It should report one stable baseline run and one stable abstain-aware run under
matched training settings.

\begin{table*}[t]
  \centering
  \caption{Main Stage 1 comparison between the answer-only baseline and the
  abstain-aware model. Replace the placeholder values once the final runs are
  selected.}
  \label{tab:stage1-main-results}
  \small
  \setlength{\tabcolsep}{8pt}
  \renewcommand{\arraystretch}{1.1}
  \begin{tabular}{@{}lcccc@{}}
    \toprule
    Model & Answerable F1 & Decision-aware Overall F1 & Unsupported Rate & Abstain F1 \\
    \midrule
    Baseline & TBD & TBD & TBD & N/A \\
    Abstain-aware & TBD & TBD & TBD & TBD \\
    \bottomrule
  \end{tabular}
\end{table*}

The baseline is trained only on answerable examples, while the abstain-aware
variant includes explicit supervision for unanswerable cases.
A Stage 1 win should be defined conservatively:
unsupported-answer rate should decrease, answerable performance should remain
usable, and abstain quality should improve without the model collapsing into
trivial refusal behavior.

Because the abstain-aware system uses a validation-selected threshold, the main
table should be read together with a threshold sweep rather than as a single
isolated operating point.

\begin{figure}[t]
  \centering
  \fbox{\rule{0pt}{1.5in}\rule{0.95\linewidth}{0pt}}
  \caption{Placeholder for the abstain-threshold sweep.
  The final figure should plot answerable $F_1$, unsupported-answer rate, and
  abstain $F_1$ against the null-score difference threshold, with the selected
  threshold marked.}
  \label{fig:stage1-threshold-sweep}
\end{figure}

\paragraph{Fill-in summary paragraph.}
Once the stable runs are selected, this section can be finalized quickly with a
summary of the following form:
\begin{quote}
Compared with the answer-only baseline, the abstain-aware model reduced
unsupported-answer rate from \textit{[X]} to \textit{[Y]} while answerable
\(F_1\) changed from \textit{[A]} to \textit{[B]}.
Decision-aware overall \(F_1\) changed from \textit{[C]} to \textit{[D]}, and
the validation-selected abstain threshold was \textit{[T]}.
Taken together, these results suggest \textit{[one-sentence interpretation]}.
\end{quote}

\paragraph{Fill-in threshold interpretation.}
The threshold-sweep discussion can also be written from a compact template:
\begin{quote}
Figure~\ref{fig:stage1-threshold-sweep} shows the expected trade-off as the
abstain threshold changes.
As the threshold moves toward \textit{[direction]}, unsupported-answer rate
\textit{[decreases/increases]} while answerable \(F_1\)
\textit{[decreases/increases]}.
The selected threshold lies in a region where \textit{[brief justification]},
which suggests the reported gain is \textit{[robust / sensitive]} rather than a
single brittle operating point.
\end{quote}

The final version of this section should answer three questions:
\begin{enumerate}[leftmargin=1.25em]
  \item Does abstention supervision reduce unsupported answers?
  \item How much answer quality is preserved on answerable examples?
  \item Is the observed gain large enough to justify the extra abstain-aware
  machinery?
\end{enumerate}

The final write-up should also include:
\begin{itemize}[leftmargin=1.25em]
  \item one short paragraph interpreting the threshold sweep,
  \item a compact qualitative comparison of representative failure cases, and
  \item a short note on whether the gain appears to come from better abstention
  behavior rather than threshold gaming.
\end{itemize}
